\documentclass[12pt]{article}
\usepackage{graphicx} % Required for inserting images
\usepackage{setspace}
\usepackage{fancyhdr}
\usepackage{amsmath}
\usepackage{caption}
\usepackage{graphicx}
\usepackage{booktabs}
\usepackage{adjustbox}
\usepackage{url}
\usepackage[utf8]{inputenc}
\usepackage{csquotes}
\usepackage{natbib}
\usepackage{pdfpages}
\usepackage[raggedright]{ragged2e}
\usepackage[hidelinks]{hyperref}
\usepackage[letterpaper,margin=1in]{geometry}
\bibliographystyle{apsr}   % or plainnat, chicago, apa, etc.
\setcitestyle{authoryear,open={(},close={)},aysep={ }}

\title{We Don't Do That: National Identity and Nuclear Non-Use} 
\author{Wyatt King}
\date{Draft Sections: February 6}

\begin{document}

\maketitle

\blockquote{“I have the greatest pride in India and in the many things that India has given to humanity, and I think those are things of the greatest value to humanity and humanity will yet profit by them. I do believe, I know my people to some extent, liking them enormously. I do know their failings too, and I do not want to make them profess one thing and do something entirely opposite to it. That is hypocrisy and cowardice. There is a grave danger. So while I am convinced of the virtue of non-violence and its power, I am not sure that people in this country, or for that matter of any other country, at the present moment are capable of carrying the burden of non-violence. And if they fail, they fail utterly.” }

\begin{center}
    - Jawaharlal Nehru speaking at the Anti-Nuclear Arms Convention, 1962
\end{center}

\section{Introduction}
In his Nobel Prize lecture, Thomas \citet{schelling_astonishing_2006} identified “the most spectacular event” of the late twentieth century as one that did not occur. “We have enjoyed sixty years without nuclear weapons exploded in anger,” he said, “Can we make it through another half dozen decades?”

At the time Schelling gave his speech, burgeoning scholarship argued that a taboo had developed around nuclear weapons during the Cold War, making their use “unthinkable” by its end \citep{tannenwald_stigmatizing_2005}. While scholars of the taboo primarily couched their argument in elite-level case studies, they extended it to public opinion, suggesting that it too exhibited a widespread aversion to the use of nuclear weapons. Over the last ten years, a new generation of scholarship on the taboo has challenged this view by fielding survey experiments in which respondents are asked how strongly they would support the first-use of nuclear weapons in different hypothetical scenarios. Consistently, these surveys have found that a majority of respondents in Western, nuclear-armed states would support a first-strike nuclear attack \citep{sagan_revisiting_2017, haworth_what_2019,allison_under_2022, dill_kettles_2022}.

Although these surveys provide robust evidence that the public's aversion to the use of nuclear weapons is not absolute, they do not dismiss that there is a norm against it. All else being equal, people prefer using conventional weapons to nuclear ones and sizable portions of public opinion still strongly oppose using nuclear weapons even when they would be strategically useful \citep{press_atomic_2013}. Consequently, if it is true that the public possesses a normative aversion to the use of nuclear weapons, then what factors can reify this norm and expand its effect?

To answer this question, I return to initial theorizing on the nuclear taboo and argue that people's conceptions of their state's national identity may be one such factor. Following from constructivist theory, I assume that states pursue patterns of behavior that cohere with their self-perceived identities \citep{tannenwald_nuclear_2007,wendt_social_1999}. Crucially, these identities can be constructed such that they disqualify certain actions as being beyond the pale--so opposed to a state's national identity as to be antithetical to its interests.

To test whether norms have such “constitutive effects,” I field an online survey experiment of Indian citizens similar to those pioneered by Scott Sagan and Benjamin Valentino. Exposing respondents to a hypothetical scenario in which India launches a preemptive counterforce attack against Pakistan, I test how national identity-based arguments opposing nuclear use change the degree of public approval for the attack. Furthermore, considering that the institutionalization of norms may affirm people's conceptions of national identity as being opposed to nuclear use, I gauge how informing respondents of India's No First Use policy changes public support for nuclear use both independent of national identity based arguments and in conjunction with them.

Ultimately, the purpose of this thesis is twofold. Narrowly, it contributes to the upswell of survey experiments from the past ten years, in which scholars have sought to understand what factors impact public support for the use of nuclear weapons. It does so by interrogating the effects that national identity and institutionalization can have, and whether those effects are significant in the manner that scholars of the taboo suggest. More broadly, it seeks to evaluate claims in constructivist literature that norms can have a “constitutive effect,“ in which people's foreign policy preferences are impacted by their impression of state identity.

[DISCUSSION OF RESULTS]

\section{Literature Review}

In political behavior literature, conceptions of national identity have oft been thought to overlap with people's policy preferences. [LOOK UP HOW NATIONAL IDENTITY HAS MANIFESTED IN IMMIGRATION POLICY AND ON OTHER TOPICS]

In international relations, the role that norms and national identity play in impacting state behavior has been marginalized by realist scholars who assume that states are strictly security seekers. Unmoved by ideational concerns, realists argue that norms only exist as a product of states' material interests. For example, the norm against violating another state's sovereignty, realists argue, is only salient when a state does not have either the material ability or the interest in conquering another state. However, once both of those conditions have been met, then the norm against violating a states sovereignty is

The “constructivist turn” brought national identity to the forefront of understanding state behavior. Whereas most realist scholars assume that states are rational security-seekers unmoved by ideational concerns, constructivism breaks down this assumption and contends that states may be influenced by norms. 


, which they define as “a standard of appropriate behavior for actors with a given identity,” (Finnemore and Sikkink 1998, 891).  

norms do not manifest exclusively through utilitarian calculations of what is in states' security interests. Rather, norms can have both regulative and constitutive effects \citep{tannenwald_nuclear_2007, wendt_social_1999}. The former refer to how norms ``constrain exogenously given self interest and behavior, or lead to recalculations by changing the ‘price’ of behavior,'' \citep[~45]{tannenwald_nuclear_2007}. In this formulation, consistent with the logic of consequences, norms are the product of states' material interests. States uphold a given pattern of behavior not because doing so follows from other-regarding moral reasoning, but because they risk facing material penalties for breaking norms.

Norms' constitutive effects, by contrast, refer to how states uphold a given pattern of behavior as a means to validate a social identity. In this interpretation, states may incorporate non-egoistic moral reasoning into their decision making by reconciling a morally dubious action with their conception of national identity. If an action is so reprehensible as to violate people's understanding of state identity, then they will refrain from performing it. In simplified terms, arguments about state identity and morality amount to saying “\textit{we} don't do \textit{that}.” Because of some deeply rooted understanding of self-hood and what conduct follows from that conception of self, certain behaviors are considered anathema to a state's interests.

In constructivist scholarship, identities are often treated as inter-subjective and shared between like-minded states, producing a similar pattern of behavior among states who claim that identity (Finnemore and Sikkink 1998). States are incentivized to comply with norms because if they do not, they will experience embarrassment, shame, anxiety, or guilt among their peers. As well as being social, norms can also be personal (Fearon 1999). Although states may be expected to act in a given way if they are “Christian,” “Marxist,” or “post-colonial,” these titles convey little about the cultural, economic, and historical conditions that make them unique relative to other states with whom they share group identities. These personal identities create another opportunity for norms to influence state behavior. Even if a given action does not violate a norm shared by an over-arching group, it can violate an aspect of one state's identity from which it derives pride, a feeling of righteousness, or a similarly salient affect. 

In the context of nuclear use, \citet{tannenwald_nuclear_2007} argues that American political and military elites opposed nuclear use by the end of the Cold War because of how it would violate both its group and personal identity. Regarding the former, Tannenwald points to how policymakers and newspaper editorial boards declaimed nuclear use because it was “uncivilized.” It was the type of action that Iraq or another “rogue state” would do. At the same time, members of the H.W. Bush administration dismissed the prospect of using nuclear weapons because it violated an unspecified aspect of American identity. Then-Chairman of the Joint Chiefs of Staff Colin Powell called using nuclear weapons ``un-American'' and even tore up a war plan describing the U.S.'s nuclear options, saying that he did not want to “let the Genie out of the bottle.” Likewise, White House Chief of Staff John Sununu said that nuclear use was something, ``we just don't do.''

What is the nuclear taboo and how does it operate? 

The nuclear taboo can broadly be understood as ``the de facto prohibition against the first use of nuclear weapons'' \citep[10]{tannenwald_nuclear_2007}. This prohibition stems from a widespread, moral revulsion to the use of nuclear weapons, which developed during the Cold War and was fully consolidated by its end. A taboo is different from a norm because breaching it is unthinkable and constitutes crossing a bright red line. Breaking the nuclear taboo is not only something that states avoid doing, as they might eschew deliberately killing civilians, but it is even unimaginable.

Although scholars of the taboo primarily relied on elite-level case studies in making their argument, they extended it to the public, suggesting that it too displays an aversion to nuclear use. \citet{tannenwald_nuclear_2007}  highlights how surveys asking respondents to reflect on the U.S.'s bombings of Hiroshima and Nagasaki have shown a steady decline in approval since the 1940s.\footnote{See \citet{sagan_revisiting_2017} for a condensed overview of this polling.} Moreover, polling conducted concurrently with various conflicts suggest that there was broad public opposition to nuclear weapons use. Before the first Gulf War, for example, more than 75\% of respondents opposed using nuclear weapons should the war have become a stalemate, and a slim majority opposed the use of nuclear weapons even if Iraq attacked American troops with chemical or biological weapons (CNN/Time Magazine 1990).

Over the last ten years, a second generation of taboo research has been published, challenging whether such an aversion to nuclear use exists today. Focusing on methodological innovations instead of theoretical ones \citep{smetana_forum_2021}, new research has employed survey experiments to evaluate the extent of public support for a first strike nuclear attack in various hypothetical scenarios. \citet{sagan_revisiting_2017} notably simulated what they argue is a modern-day equivalent of the U.S.'s bombings of Japan. Contending that polling Americans about their retrospective preferences for nuclear use during World War II does not adequately prime them to think about the trade-offs of restraint, they presented respondents with a scenario in which the U.S. dropped an atomic bomb on Iran's second largest city to end a ground war that would otherwise kill 20,000 American troops. Despite respondents being told such an attack would kill two million Iranian civilians, 60\% of people approved of it.

Following the publication of these findings, similar studies have since been released affirming that large swaths—if not outright majorities—of the Western public would approve of the first-use of nuclear weapons in different hypothetical scenarios \citep{aronow_note_2019,haworth_what_2019, dill_kettles_2022,schwartz_when_2024}. As these studies have proliferated, the scope of debate has expanded beyond whether there is a nuclear taboo to what factors impact how strongly the public supports nuclear use. Broadly, scholars have focused on three considerations in answering this question: The military and strategic utility of a nuclear attack, states' legal obligations under international law, and the number of civilian casualties incurred by an attack.

\textit{Military and Strategic Utility} Among the first considerations that Sagan and Valentino entertained in their on-going effort to ascertain what factors affect public preferences for nuclear use was the military and strategic utility of an attack. Together with Daryl Press, they argue that the public overwhelmingly—though not exclusively—follows a utilitarian logic of consequences in weighing whether to support nuclear use \citep{press_atomic_2013}. According to the logic of consequences, people's support for nuclear use is predicated on their perceptions of the combined short- and long-term advantages of nuclear weapons as compared to conventional ones. By the military utility of an attack, Press, Sagan, and Valentino refer to the immediate battlefield advantages conferred by atomic weapons use. Should a nuclear bomb destroy an enemy or their infrastructure more substantially, with greater certainty, and while better mitigating domestic casualties than conventional arms, then people will be more inclined to support their use.

By the strategic utility of an attack, the authors refer to the long-run disadvantages tied to setting a new precedent for nuclear use. As Paul (2009) argues, nuclear non-use can be understood as an iterated prisoner's dilemma. On any one battlefield, it may be in a state's interests to defect and use nuclear weapons, but doing so incentivizes an adversary to do the same in a future conflict. As a result, nuclear use may not be desirable because of how it changes opponents' utility functions during future crises. According to this formulation, the apparent norm against nuclear use does not stem from moral reasoning but is the product of coordination, either implicit or explicit, between nuclear-armed states to achieve a mutually beneficial outcome.

Survey experiments have consistently shown that military and strategic calculations exert a powerful influence on the public's willingness to support nuclear weapons use. When a state has a nuclear-armed opponent, the public is far less willing to use nuclear weapons than in asymmetric dyads \citep{haworth_what_2019,allison_under_2022, smetana_moscow_2022}. Holding the opponent constant, support for nuclear use increases dramatically with the chances of successful attack relative to conventional weapons \citep{press_atomic_2013}. To test whether the strategic utility of an attack impacts public preferences for nuclear use, \citet{sagan_atomic_2024} replicated their\textit{ Hiroshima in Iran} scenario, but informed respondents that members of the Joint Chiefs of Staff worried that using weapons would set a bad precedent that could ultimately harm American security. In this treatment group, approval of the attack declined by nearly fifteen percentage points. Moreover, this reduction largely held after respondents were informed that other members of the JCS disagreed with this assessment.

\textit{International Law}
\noindent
Beyond the military and strategic utility of a nuclear attack, scholars have also considered how international law may limit states' willingness to use nuclear weapons. According to \citet{tannenwald_nuclear_2007}, the institutionalization of nuclear non-use through international agreements has cemented the taboo. Tannenwald interprets institutionalization broadly, saying that it spans not only international human rights law (IHL) but also the Nuclear Non-Proliferation Treaty, bilateral anti-ballistic missile agreements, and declared No-First Use policies. In survey experiments, however, scholars have primarily focused on how IHL can stem support for norm-breaking behavior.

According to the logic of consequences, international law may impact public opinion if people believe that, by breaking it, their state will face material repercussions for their actions. Because international law lacks an enforcement mechanism, though, \citet{sagan_does_2020} argue that it likely has little effect on public attitudes toward nuclear use. Others, however, have argued that international law may have an effect on state behavior if states follow a logic of appropriateness rather than one of consequence. According to the former, states are rule followers who, though not obligated to behave in a given manner, do so anyways to validate a social identity that they have constructed (Fearon and Wendt 2002). According to this thinking, states are not strictly security seekers as neorealists suggest. Rather, they are bound by an informal set of ideas and norms for how states similar to them are expected to behave. 

Outside of nuclear use, surveys have found that the American public is reluctant to support actions knowing that they violate international law \citep{chilton_laws_2014, kreps_flying_2014}. \citet{carpenter_stopping_2020}, replicating a scenario from \textit{Hiroshima in Iran} in which the U.S. uses conventional weapons in a saturation bombing, find that respondents' support for the bombing declines after they are informed that it violates international law. \citet{sagan_does_2020} question the relevance of their finding, however, by pointing out that the effect size of informing people about international law is relatively small and easily reversible. Replicating Carpenter and Montgomery's research, they found that while a minority of respondents approved of the saturation bombing after being informed of it violating international law, only slightly less than 50\% actually did. In a follow up looking at nuclear use, \citet{sagan_atomic_2024} further note that if there was disagreement among the Joint Chiefs of Staff regarding the legality of a strike, the initial effect of telling people about international law washed out.


\textit{Civilian Casualties and Moral Frameworks}
\noindent
Scholars have considered how the number of civilian casualties may change the extent of support for nuclear use. If it is true that states have internalized a sensitivity to excessive civilian casualties, owing in part to the legacy of aerial bombings during World War II \citep{pinker_better_2011, thomas_ethics_2014}, then it should follow that people would be opposed to a nuclear attack that would cause large-scale, unnecessary violence. On theoretical grounds, \citet{downes_desperate_2006} and \citet{sagan_revisiting_2017} dispute this characterization by arguing that if such a norm exists at all, it is not strong enough to withstand the domestic political pressure that would arise from a costly war. Downes even concludes that democratically-elected governments are more likely to commit violence against civilians because they may be sensitive to the costs of conflict and, thus, more likely to resort to desperate means. Sagan and Valentino, focusing more narrowly on public opinion, make a similar point, arguing that voters are willing to tolerate large-scale violence perpetrated against civilians in the event substantial troop casualties would otherwise occur.

To test whether civilian immunity exerts a constraining influence on nuclear use, \citet{sagan_revisiting_2017} manipulated the number of civilian casualties incurred by a first-strike attack, holding the number of domestic casualties in a would-be ground war constant at twenty thousand. Despite ratcheting the number of expected civilian deaths from one hundred thousand to two million, public approval of the attack decreased by only a fraction of a percentage point. Meanwhile, preference for the nuclear attack over continuing the ground war declined by eight percentage points, a margin that was not statistically significant in their sample.

In response to this finding, \citet{rathbun_greater_2019} argue that there are a diversity of moral frameworks that people may use in evaluating the merits of nuclear use and not strictly a logic of consequence or appropriateness as Press, Sagan, and Valentino initially argued. Using the typology of moral foundations developed by Graham et al. (2011), they find that people's support for nuclear use changes depending on whether they adopt liberal, non-liberal, or retributive moral frameworks. Based on this typology, Rathbun and Stein subset survey respondents based on whether they had “individualizing foundations,” a sensitivity to others' suffering and a concern for equality and fairness, or “binding foundations,” which emphasize deference to authority, in-group loyalty, and the subordination of individual rights to the community. Among those who had stronger individualizing beliefs, support for nuclear decreased substantially as civilian causalities rose. That being said, few respondents overall preferred using nuclear weapons to conventional ones as the number of civilian casualties approached one million, in contrast to Sagan and Valentino's earlier findings.

\section{Theory}

Even as evidence accumulates that there is no taboo against nuclear use among the public, there is reason to believe that a norm against it exists in some capacity. \citet{press_atomic_2013} found, for example, that all else being equal, people vastly prefer the use of conventional weapons to nuclear ones. Moreover, extant studies have often stacked the deck against non-use, creating “tough tests” for the taboo in which nuclear use would achieve a compelling military goal. In doing so, they have shown both that there is apparently no absolute prohibition against nuclear use but also that nearly half of the public disapproves of nuclear use even in the most favorable circumstances for their use. By consequence, it is worth considering how this norm against nuclear use can be strengthened.

\citet{sagan_atomic_2024} made a first effort to understand what arguments and counter-arguments most strongly affect public support for nuclear use, but they exclusively focused on the military utility, precedence, and legality of an attack. Although these factors are worthy of study, the first two, at least, follow from a strict logic of consequences, treating the public as security-seekers who will uphold a norm when it benefits their short- and long-term interests to do so.\footnote{As noted previously, international law does not necessarily follow from a logic of consequences, given international law's lack of a compelling enforcement mechanism. That being said, it may also not follow from a clear logic of appropriateness if the public believes that the state will be punished for violating behavior.} This focus has produced evidence to suggest that the public is most sensitive to the strategic considerations of atomic weapons use, but leaves open the question of whether the public is also motivated by ideational factors that limit people's impression of acceptable behavior.

The latter is of academic interest because according to constructivist scholarship, norms do not manifest exclusively through utilitarian calculations of what is in states' security interests. Rather, norms can have both regulative and constitutive effects \citep{tannenwald_nuclear_2007, wendt_social_1999}. The former refer to how norms ``constrain exogenously given self interest and behavior, or lead to recalculations by changing the ‘price’ of behavior,'' \citep[~45]{tannenwald_nuclear_2007}. In this formulation, consistent with the logic of consequences, norms are the product of states' material interests. States uphold a given pattern of behavior not because doing so follows from other-regarding moral reasoning, but because they risk facing material penalties for breaking norms.

Norms' constitutive effects, by contrast, refer to how states uphold a given pattern of behavior as a means to validate a social identity. In this interpretation, states may incorporate non-egoistic moral reasoning into their decision making by reconciling a morally dubious action with their conception of national identity. If an action is so reprehensible as to violate people's understanding of state identity, then they will refrain from performing it. In simplified terms, arguments about state identity and morality amount to saying “\textit{we} don't do \textit{that}.” Because of some deeply rooted understanding of self-hood and what conduct follows from that conception of self, certain behaviors are considered anathema to a state's interests.

In constructivist scholarship, identities are often treated as inter-subjective and shared between like-minded states, producing a similar pattern of behavior among states who claim that identity (Finnemore and Sikkink 1998). States are incentivized to comply with norms because if they do not, they will experience embarrassment, shame, anxiety, or guilt among their peers. As well as being social, norms can also be personal (Fearon 1999). Although states may be expected to act in a given way if they are “Christian,” “Marxist,” or “post-colonial,” these titles convey little about the cultural, economic, and historical conditions that make them unique relative to other states with whom they share group identities. These personal identities create another opportunity for norms to influence state behavior. Even if a given action does not violate a norm shared by an over-arching group, it can violate an aspect of one state's identity from which it derives pride, a feeling of righteousness, or a similarly salient affect. 

In the context of nuclear use, \citet{tannenwald_nuclear_2007} argues that American political and military elites opposed nuclear use by the end of the Cold War because of how it would violate both its group and personal identity. Regarding the former, Tannenwald points to how policymakers and newspaper editorial boards declaimed nuclear use because it was “uncivilized.” It was the type of action that Iraq or another “rogue state” would do. At the same time, members of the H.W. Bush administration dismissed the prospect of using nuclear weapons because it violated an unspecified aspect of American identity. Then-Chairman of the Joint Chiefs of Staff Colin Powell called using nuclear weapons ``un-American'' and even tore up a war plan describing the U.S.'s nuclear options, saying that he did not want to “let the Genie out of the bottle.” Likewise, White House Chief of Staff John Sununu said that nuclear use was something, ``we just don't do.''

If it is true that states' conceptions of their identity can limit what type of weapons they are willing to use, then it follows that:

\begin{quote}
\textbf{H1:} Public support for the use of nuclear weapons will decrease as people are exposed to arguments that the use of nuclear weapons is antithetical to their state's national identity.
\end{quote}

Beyond the role that identity plays in shaping state behavior, constructivist scholarship argues that the institutionalization of norms reifies them, potentially leading to their widespread internalization (Finnemore and Sikkink 1998). According to Tannenwald, what counts as the institutionalization of the nuclear taboo is diverse, spanning formal disarmament and non-proliferation treaties, testing bans that stigmatize nuclear weapons development, and anti-ballistic missile treaties that legitimize deterrence. 

As mentioned, a handful of survey experiments have looked into how priming survey respondents about international law impacts their support for nuclear use \citep{carpenter_stopping_2020, sagan_does_2020, sagan_atomic_2024}.  Focusing on international law is worthwhile and contributes to scholarly debates outside of nuclear weapons strategy regarding its effectiveness in regulating state behavior; however, pursuant Tannenwald's initial theorizing, it is worth considering how other types of institutions affect support for nuclear use. 

I propose focusing on No First Use policies, specifically. NFUs are unique among types of institutionalized non-use because they are explicit and unilateral. Compared to international law, which does not explicitly ban the use of small-yield tactical nuclear weapons or similarly proportionate nuclear attacks, NFUs do, amounting to an almost categorical prohibition on the first-use of nuclear weapons.\footnote{ NFUs may permit nuclear retaliation against chemical and biological weapons or conducting a first-strike pre-emptive attack in light of an adversary's imminent nuclear attack.} As it relates to identity, NFUs are unique in that they are unilaterally declared. Rather than being the product of international agreement, NFUs are exclusively held by one state. 

Although states will often adopt an NFU as a way to reduce the potential of a pre-emptive nuclear attack by an adversary, an NFU may also have the knock-on effect of positioning a state's identity as being antithetical to nuclear use. In states that have adopted such policies, politicians will often curate an image of being a “responsible” nuclear power (in the following section, I will discuss the creation of this identity in the Indian case). 

\begin{quote}
\textbf{H2:} Informing people that their state has a No First Use policy will decrease support for a first-strike nuclear attack.
\end{quote}

It may also be reasonable to believe that there is an interactive effect between a state's no-first use policy and their national identity. If it is true, as constructivist scholars contend, that states seek to validate social identities within the international system, then declaring a No First Use policy should be one way through which a state can affirm its identity as a responsible nuclear power. The combined effect therefore of informing survey respondents about an NFU and exposing them to pacifistic national identity arguments may be greater than one treatment on its own. As such, I propose the following hypothesis:

\begin{quote}
\textbf{H3:} Informing individuals about their state’s No First Use policy in combination with exposure to pacifistic national identity arguments will reduce support for the use of nuclear weapons more than either treatment alone.
\end{quote}

\section{Research Design}
To test whether national identity can alter support for nuclear first use, I field a survey experiment modeled off of those created by Sagan and Valentino. However, instead of fielding it in a Western nuclear armed state, as most survey experiments up to now have done, I conduct it in India. I do so because unlike Western, nuclear-armed states, India has a lasting tradition of political elites arguing that non-violence is politically pragmatic, morally obligatory, and part of Indian national identity. None did so as clearly as nationalist leader Mahatma Gandhi. According to Gandhi, people were spiritually required to uphold \textit{ahimsa} (a religiously-inflected understanding of non-violence shared by Hindus, Jains, and Buddhists that literally means “to not kill or injure”). Gandhi's understanding of \textit{ahimsa} was deontological and made no meaningful distinction between offensive and defensive wars. Following from his categorical aversion to violence, Gandhi's opposition to nuclear use was staunch and unequivocal.

After Gandhi's death, his peers in the Indian nationalist movement adapted non-violence to Indian foreign policy, at least rhetorically. In 1954, India's Prime Minister Jawaharlal Nehru became the first world leader to call for disarmament. During negotiations over what would become the Nuclear Non-Proliferation Treaty, India sought to give teeth to these demands. Criticizing nuclear-armed powers for their hypocrisy in wanting to stop proliferation while allowing themselves to maintain and enlarge their nuclear arsenals, India refused to sign a multilateral non-proliferation treaty unless nuclear weapons states were required to disarm. Otherwise, the NPT amounted to what Indian policymakers called “nuclear apartheid,” in which wealthy states would secure their access to nuclear weapons and protect a crucial strategic advantage over poorer, primarily post-colonial, states. On these grounds, India publicly rejected the NPT.

To be sure, for however much India's opposition to the NPT stemmed from sincere moral feeling, it was also cut with an understanding of strategic necessity. In 1964, China conducted its first nuclear tests and proceeded to build up its arsenal. Due to long-standing border disputes between the two countries, which culminated in the 1962 Sino-Indian War, China's possession of nuclear weapons risked tilting the balance of power in the Himalayas in its favor. By not signing the NPT, India guaranteed itself the option to later balance against China by testing nuclear weapons—as it would do in 1974. 

Even as India pursued nuclear weapons, it continued to frame itself as a non-violent and responsible power. The codename for India's nuclear test was “smiling Buddha,” and in announcing its success, Indira Gandhi deemed it a “peaceful nuclear explosion.” Aside from naming conventions that could amount to cheap talk, India did not press its nuclear advantage over Pakistan after testing. It did not invest in ballistic missile systems or submarines that would make its deterrent credible, and it did not test nuclear weapons again until more than two decades later. Moreover, India continued its advocacy for global disarmament. Indira Gandhi coordinated a group of five other states to advocate for disarmament, and in 1988, her son, Prime Minister Rajiv Gandhi, proposed a plan at a special U.N. summit on disarmament that would eliminate nuclear weapons globally by 2010. 

Compared to Congress-led governments of the mid- to late-twentieth century, the Bharatiya Janata Party has been far less averse to nuclear weapons. With the BJP leading the government for the first time, Prime Minister Atal Vajpayee authorized India's second nuclear weapons tests in 1998, despite security analysts widely believing that doing so unnecessarily inflamed tensions in South Asia. After Pakistan tested nuclear weapons in response, the two countries briefly went to war in 1999. In the wake of the Kargil War, Vajpayee announced that India would follow a No First Use policy, in which it would refrain from using nuclear weapons unless first attacked by biological, chemical, or nuclear weapons. 

Today, there is widespread doubt regarding whether India would actually adhere to its NFU during a crisis with Pakistan \citep{clary_indias_2019}. In its 2014 election manifesto, the BJP said that India's NFU should be “revised and updated.” When asked by Reuters journalists what this revision entailed, anonymous BJP sources said cryptically to “read the manifesto.” After public backlash, Narendra Modi clarified that India would maintain its NFU, but it has continued to take action that raise doubts regarding the legitimacy of these claims. In 2019, after a terrorist attack in India-controlled Kashmir killed \_\_ civilians, India struck non-disputed Pakistani territory for the first time since 1971. Then-U.S. Secretary of State Mike Pompeo claimed in his memoir that had the U.S. not been involved, the conflagration risked spiraling and leading to nuclear use. Several months after the attack Indian Defense Minister, standing at the site of India's second nuclear tests, announced that although India had previously “strictly adhered” to its NFU, “what happens in future depends on the circumstances.”

As Indian rhetoric regarding the future of its NFU has changed, it has invested in the material capacity required to launch a first-strike attack. Building out intelligence gathering operations, improving the range of its ballistic missile system, and developing multiple targetable re-entry vehicles, India has apparently sought the ability to launch a first-strike nuclear attack that would disable Pakistan's nuclear threat.

Even as Indian public officials have simultaneously raised doubt about the legitimacy of its NFU policy, they have still publicly clung to the language of India being a “responsible” nuclear state. In the same speech where he questioned that India would abide by its NFU in “future circumstances,” Singh said that “India attaining the status of a responsible nuclear nation is a matter of national pride.” After public outcry over the BJP's election manifesto, Modi called India's NFU “a reflection of cultural inheritance.” The contradictory nature of these statements suggests that although Indian elites may be disillusioned with the utility of India's NFU, they still believe that the language of national identity has valence when it comes to conveying the limited intentions of Indian foreign policy.

However, that India has contested interpretations of national identity is not strictly a problem for the sake of this study. Rather, it gives the survey experiment leverage. If respondents were to believe uniformly that Indian national identity prohibited the use of nuclear weapons, then the treatment of exposing respondents to national identity-based arguments would likely have no effect because people would have already in\\



in what Prime Minister Indira Gandhi described as a “peaceful nuclear explosion.” Negotiations of the NPT and India's subsequent testing of nuclear weapons embodied a central contradiction within Indian nuclear policy under the Congress-led governments of the mid- to late- twentieth century: That so much as it was outwardly ideational, it was also responsive to changes in the balance of power in South Asia. Foreign policy was simultaneously filtered through a conception of India as being the product of non-violence and the strategic necessity of protecting its borders.

More recently, Indian political elites have challenged both Gandhi's categorical aversion to nuclear weapons and his successors' more prudential distaste for them, emphasizing the urgency of Indian security. In response to intimations and intelligence reports suggesting that Pakistan had developed its own nuclear weapons throughout the 1980s and '90s, Prime Minister Atal Vajpayee authorized India's second nuclear tests in 1998, despite scholars widely believing that doing so only made the subcontinent more unstable and contributed to the outbreak of the Kargil War. India subsequently built up its nuclear arsenal, seeking a credible minimum deterrent that would successfully stop both Pakistan and China from threatening nuclear use and seizing Indian territory. As part of its credible minimum deterrent, India adopted a No First Use policy as to limit the risk that either state would potentially attack India in a pre-emptive nuclear attack. Today, India is the only democracy with an NFU, and only one of two having them, with the other being China. \\

After not testing nuclear weapons for 25 years following the peaceful nuclear explosion, India tested nuclear weapons in 1998 under the BJP-led government of Atal Vajpayee. Despite security scholars widely believing that \\

To test whether identity-based arguments can diminish public support for nuclear first use, we should look for the case where they are most likely to have an effect. If they lack an effect in this most likely case, then it is doubtful that they will have one in all other circumstances where they are not expected to be as salient. Among nuclear-armed states, India makes for the most\\


Among nuclear-armed states, India makes for this most likely case because unlike Western nuclear-armed states, India has a lasting tradition of elites adopting non-violence as both a political strategy and moral obligation. Moreover, leaders of India's dominant Congress party have applied .\\

On one hand, I could do so if the public does not care about identity in debates over nuclear use. If people are strict security-seeks, as structural realists argue, then the public should not be moved by arguments about what ought to be done. In this circumstance, national identity-based arguments will not affect preferences for nuclear use because they are of no material consequence. In the event, however, that the public does care about ideational concerns, I may still reject the null if people have already incorporated identity into their decision making. If people already possess a categorical aversion to nuclear use because it is something “we just don't do,” then further arguments to that effect will do little to change the degree of public opposition. As a result, the most compelling test case for my hypothesis would satisfy two somewhat contradictory conditions: 1). There should be the reasonable expectation that the public is not strictly security-seeking and may be sensitive to ideational concerns, and 2). Conceptions of national identity should not already be heavily-weighted in favor of nuclear non-use.

\textit{Ex ante}, India satisfies these two conditions. Unlike Western, nuclear-armed states, India has a lasting tradition of elites adopting non-violence as both a political strategy and moral obligation. None did so more actively as Mahatma Gandhi. According to Gandhi, people were spiritually required to uphold \textit{ahimsa} (a religiously-inflected understanding of non-violence shared by Hindus, Jains, and Buddhists that literally means “to not kill or injure”). Gandhi's understanding of \textit{ahimsa} was deontological and made no meaningful distinction between offensive and defensive wars. Following from his categorical aversion to violence, Gandhi's opposition to nuclear use was staunch and unequivocal.

After the U.S. bombed Japan, Gandhi called for global disarmament. Reflecting on the destruction of Hiroshima, he said that, “Unless now the world adopts non-violence, it will spell certain suicide for mankind.” When challenged that nuclear weapons could be an instrument to achieving peace by creating stable deterrence, Gandhi was incisive. He compared nuclear weapons' deterrent effect to a “a man glutting himself with dainties.” He refrains from violence “only to return with redoubled zeal after the effect of nausea.”


 Even after testing nuclear weapons, however, India maintained its posture as a responsible nuclear power. \\



So much as India has this history, however, under the Bharitya Janata Party, it has adopted a more forceful foreign policy. Whether it be skirmishing with China along their disputed border, striking uncontested Pakistani territory for the first time since 1971, or assassinating a Sikh nationalist in Canada, India's recent military and covert actions reveal a state more willing to use force to guarantee its territorial integrity and religiously-inspired identity. 

but this non-violence, couched in Gandhian interpretations of Indian nationalism, has been contested by Hindu nationalist politicians who are more willing to use force to secure India's material interests.

\subsection{Scrap}

Indian nationalists explicitly incorporated non-violence into their conception of state identity during the twentieth century. Moreover, Indian leaders were among the staunchest advocates for global disarmament dating to the early 1950s, couching their rationale in Gandhian non-violence. 

So much as India has this history, however, under the Bharitya Janata Party, it has adopted a more forceful foreign policy. Whether it be skirmishing with China along their disputed border, striking uncontested Pakistani territory for the first time since 1971, or assassinating a Sikh nationalist in Canada, India's recent military and covert actions betray its legacy of non-violence. What's more is that members of the BJP have grounded India's new foreign policy in religious nationalism, such as when Foreign Minister S. Jashainker wrote a book describing what lessons Indian foreign policy can take from the \textit{Ramayana}, a Hindu epic.

Given these diverging interpretations of Indian national identity and its interaction with foreign policy, why does India still make for a compelling test case? 

They may not have done so because people are fundamentally unconcerned about national identity, in which case my first hypothesis would always be rejected. However, it may also happen because 

They may not have done so because people are fundamentally unconcerned about national identity, or because they have already 

Among these interpretations, Gandhianism is the most skeptical of nuclear weapons. During the nationalist movement, Mohandas Gandhi called for a near-absolute adherence to non-violence amid the popular struggle against British imperialism. His opposition to  violence stemmed from an ascetic interpretation of \textit{dharma} that was not widely shared in the early twentieth century began but one that he successfully politicized. According to Gandhi, people were morally and spiritually required to uphold \textit{ahimsa} (a religiously-inflected understanding of non-violence shared by Hindus, Jains, and Buddhists that literally means “to not kill or injure”). Notably, Gandhi's interpretation of \textit{ahimsa} was deontological and made no meaningful distinction between offensive and defensive wars. Following from his categorical aversion to violence, Gandhi's opposition to nuclear use was staunch and unequivocal.

Soon after the U.S. bombings of Japan, Gandhi called for global disarmament. Reflecting on the destruction of Hiroshima, Gandhi said, “Unless now the world adopts non-violence, it will spell certain suicide for mankind.” When questioned whether nuclear weapons could achieve peace through deterrence, Gandhi was equally incisive, saying that nuclear weapons' deterrent effect would be like, \blockquote{A man glutting himself with dainties to the point of nausea and turning away from them only to return with redoubled zeal after the effect of nausea.}

In international debates over proliferation after Gandhi's death, India became among the most high-profile states calling for global disarmament, with its opposition often couched in the language of non-violence. In a speech before Parliament in 1954, Jawaharlal Nehru became the first world leader to advocate for disarmament, advocating for a “standstill” agreement in which all powers would halt nuclear weapons testing and reduce the size of their arsenals. India sought to give teeth to these demands during negotiations over the Nuclear Non-Proliferation Treaty. Criticizing nuclear-armed powers for their hypocrisy in seeking to stop proliferation while allowing themselves to maintain nuclear arsenals, India argued that the NPT amounted to “nuclear apartheid,” in which wealthy states would secure their access to nuclear weapons and protect a crucial strategic advantage over poorer, primarily post-colonial, states. On these grounds, India publicly refused to sign the NPT.

So much as Indian leaders were  inspired by normative concerns, they took a more prudential approach to violence than Gandhi's categorical aversion to it. India's first prime minister Jawaharalal Nehru understood violence as something to be avoided at all costs, but as being ingrained in state and human behavior. Speaking at the 1962 Anti-Nuclear Arms Convention in New Delhi, Nehru reflected on this contradiction, saying that non-violence could not be conflated with permissive cowardice, in which a state did not defend itself from attack because of a moral obligation to be non-violent.



\section{Scrap}

As survey experiments like those developed by Sagan and Valentino continue to be fielded, there are several empirical gaps that have yet to be filled in the literature. Among the most notable is that overwhelmingly surveys have been conducted in the United Sates, which considering its history as the only state to have used nuclear weapons may undermine efforts to draw broader conclusions regarding the status of international non-use norms. Although there is evidence suggesting that attitudes in Israel, the U.K., and France toward nuclear use are similar to those in the U.S. \citep{dill_kettles_2022}, the dearth of extant surveys conducted outside of Western, nuclear armed states makes drawing system-wide conclusions difficult. \citet{smetana_moscow_2022} found that a majority of Russian respondents would oppose the first use of nuclear weapons against a NATO member, but given NATO's mutual defense clause and its forward positioning of nuclear weapons in Europe, this conclusion is hardly surprising. As of early 2020, \citet{onderco_ideology_2022} found that sweeping majorities of the Dutch and German public opposed the first use of American nuclear weapons stationed on their territory. That being said, this opposition has substantially weakened since Russia's invasion of Ukraine \citep{onderco_hawks_2023}.

In South Asia, some surveys suggest that a majority of respondents would be willing to use nuclear weapons in certain circumstances. \citet{schwartz_when_2024}, for example, found that most Indian respondents would support using a nuclear weapon against a terrorist cell in Yemen plotting to attack India with a dirty bomb. \citet{clary_public_2021} meanwhile found that people in Punjab, Pakistan would strongly support an escalatory versus de-escalatory response to conventional Indian provocation, albeit the authors did not explicitly distinguish between whether such an escalation would be nuclear or conventional. To the author's knowledge, no survey experiment has yet to be conducted in China.


\section{Research Methods}
To test my hypotheses, I intend to employ a survey experiment in which Indian respondents are presented with a hypothetical crisis between India and Pakistan. The survey will be conducted via Amazon Mechanical Turk because of its cost efficiency relative to other convenience samples. Like other survey platforms in India, MTurk's samples are disproportionately young, male, upper-caste, and wealthy \citep{boas_recruiting_2020}. To assign treatments, I will use a randomized experiment, where respondents are assigned to one of four groups. The control group will receive no information about India's NFU and no arguments condemning nuclear use. One treatment group will receive information about India's NFU but no arguments, and another treatment will read arguments condemning nuclear use but not receive information about India's NFU. Lastly, one treatment group will both be told of India's NFU and read arguments condemning nuclear use.

Drafts of the hypothetical news articles that respondents in each treatment will read are attached in the appendix below. The situation described is consistent with India's Cold Start doctrine and inspired by what security scholars consider to be the most likely case of Indian nuclear first-use \citep{kapur_ten_2008, clary_indias_2019}. Implemented after the Kargil War, Cold Start is Indian military doctrine enabling it to quickly mobilize along the border with Pakistan following a Pakistan-backed terrorist attack. Within a few days, India can then push across the borders with Punjab and Kashmir in an effort either to claim disputed territory of force a favorable compromise in which Pakistan stops sponsoring cross-border terrorism. 

Although India publicly maintains a no first use policy, Christopher Clary and Vipin Narang argue that during such a ground-war India may contemplate launching a first strike nuclear attack. Because India's military is conventionally superior to Pakistan's, the latter may be incentivized to use a small-yield tactical nuclear weapon on the battlefield to hold territory and stop a further Indian incursion. However, given that Pakistan may use a nuclear weapon, India is in turn encouraged to preempt. Fearing retaliation against one of its cities, Clary and Narang argue that India may conduct a counter-force attack, in which it would target Pakistan's nuclear silos, airfields, and missile launchers in an effort to marginalize the chance of a nuclear attack. 

\textit{A priori}, it is difficult to know how successful such an attack would actually be, given uncertainty regarding where Pakistan's nuclear weapons facilities are, how effective its air defenses would be, and a range of other factors. In the article, I note that it is unclear whether all of Pakistan's strategic nuclear weapons have been destroyed, but that Indian government officials speaking on the condition of anonymity believe that there is a low chance of retaliation. I do so as to create a more difficult test for the taboo while not misleading respondents into believing that a counterforce attack would be guaranteed to work. Furthermore, although it is difficult to gauge how many civilians would die in a counterforce attack, I note that the U.N. warns that it is possible hundreds of thousands did. Again, I do so as to not mislead respondents into believing that a counter force attack could be committed with little harm to civilians, even if the extent of that harm is ambiguous.

Within the article, key pieces of information that may affect respondents' support for nuclear use are highlighted in the article title, dek, and pull quotes. Specifically, I emphasize the number of expected casualties, the target of the attack being Pakistan's nuclear weapon facilities, that government officials believe there is little chance of retaliation, and that India launched the attack amid a ground-war with Pakistan.

Across treatment and control groups, the article is identical, except that people assigned to learn about India's no-first use policy read an additional sentence saying that ``the nuclear attack violates long-standing Indian policy that it would never use nuclear weapons unless first attacked with weapons of mass destruction.'' A pull quote in the middle of the article is meant to draw further attention to this fact.

After reading the article, respondents in the control and NFU treatment groups immediately begin answering questions regarding their support for the use of nuclear weapons in the scenario described. People assigned to reading national identity-based arguments, however, are provided three arguments to read criticizing the use of nuclear weapons. Possible arguments that respondents will read are included in the appendix. Broadly, the arguments that I have written are meant to address different aspects of how conceptions of India's national identity impact its foreign policy. Some emphasize India's legacy of non-violence during the nationalist movement, others more explicitly address India's history of advocating for global disarmament, and a third kind emphasize a transcendent identity shared between Indians and Pakistanis. These arguments follow from previous academic literature on how India's strategic culture, national identity, and nuclear weapons policy overlap \citep{hymans_psychology_2006, bajpai_indian_2014}. To decide what survey questions to include in my experiment, I hope to field a short preliminary experiment in which respondents read the vignette and then grade how convincing a randomly assigned article is on a Likert scale. I would then include whatever arguments were graded as the most persuasive in my final experiment. That being said, considering I may lack funding, I may also forgo this trial experiment and use the first three arguments included in the appendix. I would chose these three in particular because they address different aspects of India's history vis-a-vis nuclear weapons and its identity relative to Pakistan.

When it comes to measuring how strongly respondents support the use of nuclear weapons, I intend to follow \citet{sagan_revisiting_2017} and ask respondents two questions: How strongly would they have preferred using nuclear weapons compared to continuing the ground war, and how strongly would they approve of the use of nuclear weapons? Asking respondents the former provides information about what they believe to be reasonable conduct during war. Asking them the latter is a good indication of respondents' willingness to support the prime minister's decision to use nuclear weapons. This level of approval is also politically relevant insofar as respondents should prefer continuing the ground war but still approve of the use of nuclear weapons, then the government likely faces less political risk regardless of public preferences. To further probe the political implications of nuclear use, I intend to ask respondents whether nuclear use would make them more or less likely to vote for the sitting government and whether they would publicly protest the use of nuclear weapons. These questions are meant to be tougher tests for the degree of public outcry to the use of nuclear weapons. 

Note that in keeping with Sagan and Valentino's survey experiments, I intend to use a six-point Likert scale to measure respondents' preference for and approval of nuclear use. On the scale, respondents must either approve or disapprove and cannot indicate that they are unsure or indifferent to the two options. This design mirrors the ``trolley car'' problem, which philosophers use to assess people's moral reasoning when faced with difficult ethical trade-offs \citep{sagan_does_2020}.

In addition to the above, I intend to collect information on potential covariates in a pre-treatment survey administered to respondents. I specifically ask how strongly people approve the Indian National Congress and Bharatiya Janata Party on a seven point Likert scale. These questions will serve as controls for people's political preferences in the regression models described below. I will also ask how strongly respondents agree with the statements that ``India should not use force against another country unless absolutely necessary'' and ``India should never use nuclear weapons unless it is retaliating against another country for using them first.'' Responses to these questions are meant to be highly correlated with their preferences for nuclear use after reading the hypothetical as to reduce standard errors and increase statistical power when I perform my statistical analyses.

To evaluate the support for my hypotheses, I intend to use two statistical procedures primarily. First is a difference of proportions test, comparing people's approval and preference for the nuclear attack between the control and treatment groups. Second is a logistic regression, in which I dichotomize people's approval for nuclear use as the outcome variable and include controls for people's demographics, their political beliefs, and their pre-treatment attitudes toward nuclear use. For robustness, I will also conduct a series of ordinal regressions, using the same independent variables as for the logistic regression.


\bibliography{references}

\end{document}